\documentclass[12pt,a4paper]{article}
\IfFileExists{pt-common.sty}{\usepackage{pt-common}}{\usepackage{../shared/pt-common}}
\usepackage[margin=2.5cm]{geometry}
\usepackage{setspace}
\onehalfspacing

% Prediction environment (used in §5; shares the theorem counter)
\newtheorem{prediction}[theorem]{Prediction}

\title{\textbf{Holonomy without Potential:\\
Persistence Theory and the\\
Aharonov--Bohm / Non-Locality Debate}}
\author{Yan Senez\\
\texttt{yan.senez@protonmail.com}}
\date{May 2026}

\begin{document}
\maketitle

\begin{abstract}
The foundational debate between (a) the reality of the vector
potential (Bohm--Hiley), (b) the intrinsic non-locality of the
field (strong readings of EPR/Bell), and (c) the primacy of
holonomy (Wu--Yang--Healey), admits, within the classical
field-on-manifold framework, no experimental prediction that
adjudicates between them. Persistence Theory (PT), derived
entirely from the symmetry parameter $s = 1/2$, supplies a
structural theorem and a falsifiable prediction that do so.

We show that (i) the fundamental dynamical field is unique
(the prime-gap sequence, BA0/T0); (ii) the classical potential-field
pair $A_\mu \leftrightarrow F_{\mu\nu}$ is entirely derived from
the gradient and Hessian of the scalar persistence potential
$S = -\ln \alpha$, and is therefore secondary; (iii) the primitive
object is the modular holonomy angle $\theta_p$, gauge-invariant
and topological; and (iv) the correlation $I(\Phi_p\,;\,\Phi_q\,|\,
G(\mu))$ between distinct modular channels is invariant under the
emergent geometric family $G(\mu)$ (theorem proved in the companion
paper \cite{PT_CRT_DECOUPLING}). This invariance reclassifies the
apparent spatial non-locality of Bell--EPR as an arithmetic
intra-channel CRT correlation, and the Aharonov--Bohm phase as a
topological invariant of the $\theta_3$ channel independent of the
metric. Prediction QH1a (mutual information invariant under spatial
geometry) is testable on cold-atom, trapped-ion, and diatomic
molecular platforms.

Theorems~4.1 and~4.2, together with their structural framework, are
formally verified in Lean~4 + Mathlib under the namespace
\texttt{PT.CrtDecoupling} of the companion paper
\cite{PT_CRT_DECOUPLING} (Phases~1 to~4.1, 0~\texttt{sorry},
0~\texttt{axiom}).
\end{abstract}

\tableofcontents

\section{The foundational debate: three classical positions}
\label{sec:debate}

Since the founding paper of Aharonov and Bohm
\cite{AharonovBohm1959}, quantum electrodynamics has posed a
foundational question that resists any clean experimental
resolution: which among the 4-potential $A_\mu$, the field tensor
$F_{\mu\nu}$, and the holonomy phase $\oint_\gamma A \cdot d\ell$
should be regarded as the physically real object? Three
positions divide the contemporary literature, with no classical
experiment so far able to distinguish them. This section recalls
them in their sharpest form, for the sole purpose of fixing what
Persistence Theory (PT) claims to adjudicate. The review is not
exhaustive --- for the philosophical developments we refer to the
monographs of Healey \cite{Healey2007} and Maudlin
\cite{Maudlin2011}.

\subsection{Position A: reality of the vector potential}
\label{ssec:position_A_en}

Position A, inherited from the Bohm--Hiley programme
\cite{BohmHiley1993}, holds that the potential $A_\mu$ is
physically real and that the tensor $F_{\mu\nu} = \partial_\mu A_\nu
- \partial_\nu A_\mu$ is a derived quantity. The empirical argument
is the Aharonov--Bohm effect~\cite{AharonovBohm1959}: the
interferometric phase shift $\Delta\varphi = (e/\hbar) \oint A \cdot
d\ell$ recorded by a particle that loops around an ideal solenoid
is measurable even where $E = B = 0$ along its entire trajectory.
No field crosses the path; no local object can therefore account
for the effect. The potential, by contrast, is non-zero (up to
gauge) along the trajectory, and the Bohmian programme reads this
as proof that $A_\mu$ must be regarded as an object of higher
ontological status.

The known difficulty of position A is gauge dependence. The
potential $A_\mu$ is defined only up to a transformation
$A_\mu \mapsto A_\mu + \partial_\mu \Lambda$, and every physically
measurable quantity must be invariant under this transformation.
The Aharonov--Bohm phase is gauge-invariant \emph{only as a closed
integral}; its local values are not. Position A must therefore
either accept a physical reality that is not itself gauge-invariant
--- which many consider a major ontological concession --- or
treat $A_\mu$ as a hybrid object only the gauge-equivalence class
of which is real.

\subsection{Position B: intrinsic non-locality of the field}
\label{ssec:position_B_en}

Position B accepts $F_{\mu\nu}$ as the real object but makes it
intrinsically non-local. In the strong reading proposed by Mead
\cite{Mead2000} and in the wake of the analyses of Bell
\cite{Bell1964} and the experimental tests of Aspect
\cite{Aspect1982} and Hensen et al.~\cite{Hensen2015},
the electromagnetic field is not a distribution on classical
spacetime: it carries a correlation structure between spatially
separated points that is not explained by the local variation of
its components. The Aharonov--Bohm phase then becomes the signal
of such non-locality, and the EPR correlations its extension to
matter.

The structural objection to position B is that it requires
modifying the classical ontology of the field to make it
intrinsically non-local, without specifying which mathematical
mechanism implements the modification. The available formal
machinery (connection fibre bundles, gauge theories, Wilson
operators) admits strictly local interpretations under a suitable
choice of variables. The non-locality of position B remains, in
most formulations, a property of the \emph{interpretation} of the
field rather than of its mathematical structure.

\subsection{Position C: primacy of holonomy}
\label{ssec:position_C_en}

Position C, due to Wu and Yang \cite{WuYang1975} and developed by
Healey \cite{Healey2007} and Belot \cite{Belot1998}, holds that
neither of the two local objects $A_\mu, F_{\mu\nu}$ is
fundamental: $A_\mu$ because it depends on gauge, $F_{\mu\nu}$
because it does not contain the Aharonov--Bohm phase information
in regions where it vanishes. The only object that is at once
gauge-invariant and sensitive to the Aharonov--Bohm experiment is
the Wilson loop $W_\gamma = \exp\bigl[i (e/\hbar) \oint_\gamma A
\cdot d\ell\bigr]$, a non-local functional of closed paths $\gamma$.
This object, and not $A_\mu$ or $F_{\mu\nu}$, deserves the status
of primitive.

Position C is conceptually satisfying --- it dissolves the
asymmetry between the gauge-dependence of $A_\mu$ and the
incompleteness of $F_{\mu\nu}$ --- but it leaves open the question
of the \emph{origin} of Wilson loops: why these particular
non-local objects and not others? Classical gauge theory introduces
them as a post hoc mathematical construction from a connection
that remains itself primary in the formalism. Position C therefore
remains, in a sense, a description rather than an explanation.

\subsection{The operational impasse}
\label{ssec:impasse_en}

Positions A, B, C make the same experimental predictions within
standard quantum electrodynamics. The Aharonov--Bohm phase is
measurable in all three frameworks; the EPR/Bell correlations are
reproduced in all three; cross sections and lifetimes are
identical. No measurement to date has been able to distinguish
the three positions at the level of quantitative prediction. The
debate is thus, at present, a debate of interpretation.

Persistence Theory modifies this state of affairs on two counts.
On one hand, it promotes position C to a \emph{theorem} rather
than an interpretation: the modular holonomy angle $\theta_p$ is
the unique primitive invariant satisfying four structural
conditions that neither $A_\mu$ nor $F_{\mu\nu}$ can meet
simultaneously (Section~\ref{sec:holonomy_primitive}). On the
other hand, it produces an experimental prediction --- prediction
QH1a of Section~\ref{sec:qh1a_prediction} --- whose empirical
verdict effectively distinguishes PT (and with it position C)
from a standard theory on curved manifold. The foundational
debate is no longer merely a debate of interpretation: under PT,
it becomes adjudicable.

\section{The unique field and the modular connection of PT}
\label{sec:field_connection}

Before evaluating position C in the PT framework
(Section~\ref{sec:holonomy_primitive}), three structural facts
must be established: that PT admits a unique dynamical degree of
freedom (\S~\ref{ssec:BA0_en}), that this degree of freedom
naturally carries a modular gauge connection
(\S~\ref{ssec:gauge_connection_en}), and that this connection
induces the holonomy angles $\theta_p$ through an explicit
algebraic identity (\S~\ref{ssec:T6_en}). All three facts are
established in the monograph~\cite{PT_Monograph} and reproduced
here without proof.

\subsection{The unique dynamical field BA0/T0}
\label{ssec:BA0_en}

The closure theorem T0 of chapter~25 of the
monograph~\cite{PT_Monograph} establishes that the sequence of
gaps between consecutive primes,
\begin{equation}
  g_n \;:=\; p_{n+1} - p_n \qquad (n \geq 1),
  \label{eq:g_n_def_en}
\end{equation}
is the \emph{unique} dynamical degree of freedom of the
Eratosthenes sieve, under the structural conditions U1--U4
(antidiagonality of $T_3$, GFT identity, canonical geometry of
$\DKL$ and Fisher, fixed point $\mustar = 15$). Postulate BA0 of
the monograph --- that $\{g_n\}_{n \geq 1}$ is \emph{the}
dynamical field of the theory --- is therefore promoted from
postulate to theorem, and the theory admits \emph{zero} dynamical
axioms beyond the standard tools of number theory.

This is what distinguishes PT radically from classical gauge
theories: whereas the latter must postulate a vector field $A_\mu$
as the primitive structure of electrodynamics, PT derives
\emph{all} of the gauge structure --- including the very
existence of a connection --- from the single scalar field
$\{g_n\}$ and the arithmetic coherence constraints. The gauge
connection will be, in what follows, a \emph{derived} object, not
a postulate.

\subsection{The modular gauge connection}
\label{ssec:gauge_connection_en}

For any modulus $m \geq 2$, the additive recurrence
\begin{equation}
  r_{n+1}^{(m)} \;=\; \bigl(r_n^{(m)} + g_n\bigr) \bmod m
  \label{eq:gauge_recursion_en}
\end{equation}
expresses that the trajectory of modular residues
$r_n^{(m)} = p_n \bmod m$ is entirely determined by the initial
datum and the sequence $\{g_n\}$. This recurrence is the
\emph{parallel transport equation} associated with the modular
connection of PT: at each step $n$, the residue is translated by
the ``gauge gain'' $g_n \bmod m$, and one reads the trajectory as
a transport along the ``arithmetic time'' $n$.

The gauge connection theorem~\cite[ch.~1, Thm.~gauge]{PT_Monograph}
formalises this reading: the
recurrence~\eqref{eq:gauge_recursion_en} defines, for every
squarefree integer $m = q_1 \cdots q_k$ that is a product of
distinct primes, a connection on the trivial bundle
$\Omega \times (\mathbb{Z}/m\mathbb{Z}) \to \Omega$ above the
path space $\Omega = \mathbb{N}^{\mathbb{N}}$. By the CRT
decomposition
\begin{equation}
  \mathbb{Z}/m\mathbb{Z}
  \;\cong\;
  \prod_{i=1}^{k} \mathbb{Z}/q_i\mathbb{Z},
  \label{eq:CRT_iso_en}
\end{equation}
this connection splits into $k$ independent components, one per
prime $q_i$, and one can treat each factor
$\mathbb{Z}/q_i\mathbb{Z}$ as a discrete circle of $q_i$ points
carrying its own gauge connection. It is this decomposition that
makes the modular connection of PT structurally simpler than a
classical connection on $\mathbb{R}^4$: it is, by construction,
a tensor product of finite-dimensional objects.

\subsection{The holonomy identity T6}
\label{ssec:T6_en}

The holonomy theorem T6~\cite[ch.~6]{PT_Monograph} gives the
explicit expression of the holonomy angle associated with parallel
transport along a cycle of the modular connection. For each prime
$p$ and each scale $\mu > 0$,
\begin{align}
  \delta_p(\mu) &\;:=\; \frac{1 - q(\mu)^p}{p},
  &
  \cos \theta_p(\mu) &\;:=\; 1 - \delta_p(\mu),
  \label{eq:delta_theta_en}
\\
  \sinp{p}(\mu) &\;=\; \delta_p(\mu)\bigl(2 - \delta_p(\mu)\bigr),
  &
  q(\mu) &\;\in\; \{\qplus,\; q_-\},
  \label{eq:sin2_theta_en}
\end{align}
where $\qplus(\mu) = 1 - 2/\mu$ and $q_-(\mu) = e^{-1/\mu}$ are
the two canonical values of $q$ singled out by the
maximum-entropy bifurcation L0~\cite[ch.~3]{PT_Monograph}. For the
holonomy questions treated here we adopt the coupling branch
$q = \qplus$, which describes the electromagnetic sector and the
leptons.

Identity~\eqref{eq:sin2_theta_en} is purely algebraic --- it uses
only the trigonometric definition
$\sin^2 \theta = 1 - \cos^2 \theta = 1 - (1 - \delta)^2 =
\delta (2 - \delta)$ --- but it has a striking consequence: the
usual trigonometry, and with it the constant $\pi$, \emph{emerges}
from the modular dynamics without having been introduced at the
start. The continuous circle is never postulated; it appears as
the continuous limit of the discrete circle
$\mathbb{Z}/p\mathbb{Z}$ as $p \to \infty$, and $\pi$ appears via
the Riemann zeta function $\zeta(2) = \pi^2/6 =
\prod_p (1 - 1/p^2)^{-1}$, which directly relates the measure of
the continuous circle to the sieve of primes.

\paragraph{Structural summary.}
The three elements established in this section sketch an
architecture in which the gauge connection and the holonomy
angle are \emph{derived} objects, computable from the single
scalar field $\{g_n\}$, and not independent postulates. Position
C of the foundational debate --- holonomy as primitive --- takes,
under PT, a substantially stronger sense than it has in field
theory on manifold: it is not an option of interpretation but
the unique outcome of a structural construction that admits no
other choice. This is the status that the following section
formalises.

\section{The holonomy angle $\theta_p$ as primitive object}
\label{sec:holonomy_primitive}

This section establishes the fundamental result that distinguishes
PT from classical gauge theories: under the structural hypotheses
of PT, the modular holonomy angle $\theta_p$ is the
\emph{unique} primitive invariant simultaneously satisfying four
natural conditions, and the classical pair $(A_\mu, F_{\mu\nu})$
is entirely reconstructible as a derivative of a more primitive
object --- the scalar persistence potential $S = -\ln \alpha$.
Position C of the foundational debate is no longer an
interpretation: it is a structural consequence.

\subsection{Proposition 3.1: $A_\mu$ and $F_{\mu\nu}$ are derived}
\label{ssec:prop_3_1_en}

\begin{proposition}[Persistence Potential Decomposition]
\label{prop:A_F_derived_en}
Let $S(\mu) = -\ln \alpha(\mu)$ be the persistence potential. Then
\begin{itemize}
  \item the 4-potential $A_\mu$ is identifiable with the gradient
  $\partial_\mu S$;
  \item the field tensor $F_{\mu\nu}$ is identifiable with the
  antisymmetric components of $\mathrm{Hess}(S)$;
  \item the metric $g_{\mu\nu}$ is identifiable with the symmetric
  components of $\mathrm{Hess}(S)$ (Lemma F, ch12 monograph).
\end{itemize}
\end{proposition}

\begin{proof}[Proof]
The construction is detailed in chapter~12 of the
monograph~\cite{PT_Monograph}; we recall here the three key steps.

\textbf{Step 1 (gradient identification).}
The Fisher metric associated with the stationary distribution
$\pi_m$ of the sieve is computed as the Hessian of the logarithm
of the partition function $Z_{\rm Ruelle}$, and the function
$\alpha(\mu)$ appears as the first eigenfunction of the transfer
operator. By construction, the gradient
$\partial_\mu S = -\partial_\mu \ln \alpha$ plays in the PT formalism
the role played by $A_\mu$ in classical electrodynamics: it enters
linearly in the interaction term of the effective Lagrangian
(\S~F3 of the monograph) and undergoes, under modular gauge
transformations $\theta_p \mapsto \theta_p + 2\pi k$, a phase
shift compatible with classical gauge invariance.

\textbf{Step 2 (antisymmetric Hessian identification).}
The antisymmetric components of the Hessian
$\mathrm{Hess}(S)_{\mu\nu} = \partial_\mu \partial_\nu S$ encode
the curvatures derived from the scalar potential, and reproduce
the six degrees of freedom of $F_{\mu\nu} = \partial_\mu A_\nu -
\partial_\nu A_\mu$ by direct identification. The identity
$dA = F$ becomes, under PT, the tautological identity
$d(dS) = 0$ projected onto the antisymmetric components of
$\mathrm{Hess}(S)$. This projection is non-trivial in the presence
of the Lorentzian signature $g_{00} < 0$, which separates the
spatial and temporal components of the Hessian.

\textbf{Step 3 (symmetric Hessian identification).}
The symmetric components of $\mathrm{Hess}(S)$ reproduce the
metric $g_{\mu\nu}$ of the emergent spacetime (Bianchi I on the
active primes $\{3, 5, 7\}$; see Lemma F of
ch.~12~\cite{PT_Monograph}). This identification is, in itself,
the content of the metric-reconstruction theorem: the
spacetime metric of PT \emph{is} the Hessian of the logarithm of
the partition function, with no new structure introduced.
\end{proof}

The immediate consequence is that $A_\mu, F_{\mu\nu}, g_{\mu\nu}$
are all three \emph{derivable from the same scalar potential}
$S(\mu)$, and that this potential is, in turn, computed from the
sole field $\{g_n\}$. None of these three objects is therefore
primary in the PT architecture: they are all three secondary
relative to $S$, and $S$ is itself secondary relative to
$\{g_n\}$. This structural observation underwrites the uniqueness
argument of the following subsection.

\subsection{Proposition 3.2: uniqueness of $\theta_p$ as primitive invariant}
\label{ssec:prop_3_2_en}

\begin{proposition}[Holonomic Primitiveness]
\label{prop:theta_p_unique_en}
The holonomy angle $\theta_p$ is the unique primitive invariant
simultaneously:
\begin{enumerate}[label=(\roman*),nosep]
  \item gauge-invariant,
  \item sensitive to Aharonov--Bohm,
  \item bounded,
  \item PT-definable from BA0 alone.
\end{enumerate}
\end{proposition}

\begin{proof}[Proof]
We show by exclusion that no other natural invariant
satisfies the four conditions simultaneously.

\textbf{Exclusion of $A_\mu$.}
The potential $A_\mu$ violates condition (i): it depends on gauge
by construction. No natural projection of $A_\mu$ is at once
gauge-invariant and bounded while remaining sensitive to
Aharonov--Bohm in the region where $F = 0$. Position A of the
foundational debate is therefore structurally incompatible with
the four conditions.

\textbf{Exclusion of $F_{\mu\nu}$.}
The field tensor $F_{\mu\nu}$ satisfies (i), (iii), (iv) but
violates condition (ii): it is, by construction, identically zero
in the interior of an ideal Aharonov--Bohm path, and so cannot
account for the observed phase shift. Position B of the debate
must add a supplementary non-local structure to recover sensitivity
to AB; that added structure is not PT-definable from BA0 alone.

\textbf{Exclusion of arbitrary functions of $\{g_n\}$.}
An arbitrary function $F[\{g_n\}]$ may satisfy (i), (ii), (iv) but
generally violates (iii): without a modular parallel-transport
structure, the function is not naturally bounded and does not
admit a simple identification with a measurable topological
invariant. Only functions that reduce to the modular angle
$\theta_p$ of a CRT channel $\mathbb{Z}/p\mathbb{Z}$ satisfy the
four conditions simultaneously.

\textbf{Existence and uniqueness.}
The angle $\theta_p$ defined by the
identity~\eqref{eq:sin2_theta_en} satisfies (i) (it is
gauge-invariant because defined as a closed integral),
(ii) (it is non-zero in every region where the modular connection
is non-trivial, including in the interior of an AB loop), (iii)
(it takes values in $[0, \pi]$ by construction), and (iv) (it is
defined by identity~\eqref{eq:sin2_theta_en}, which uses only
$\{g_n\}$ via $q(\mu)$). It is the unique primitive invariant
realising the four conditions.
\end{proof}

\subsection{The Wu--Yang--Healey position promoted to PT theorem}
\label{ssec:wu_yang_promoted_en}

Propositions~\ref{prop:A_F_derived_en}
and~\ref{prop:theta_p_unique_en} have a common reading: position
C of the foundational debate
(\S~\ref{ssec:position_C_en}) --- the primacy of holonomy --- is
no longer, in the PT framework, an interpretive choice among
others. It is a structural consequence: under the axioms of PT,
$\theta_p$ is the \emph{unique} primitive invariant available,
and $A_\mu, F_{\mu\nu}, g_{\mu\nu}$ are all three derived from
it via the scalar potential $S$.

This status has two immediate consequences. On one hand, PT
\emph{predicts} that any measurement of $A_\mu$ or $F_{\mu\nu}$
reduces, in its informational content, to a measurement of
$\theta_p$ on an active modular channel. No physical information
is lost by this reduction --- the AB phase is entirely encoded
in $\theta_3$ (cf.~\S~\ref{ssec:ab_consequence_en}). On the other
hand, PT \emph{excludes} the ontologies of position A (the reality
of the potential) and position B (the intrinsic non-locality of
the field): these ontologies postulate primitive objects that, in
the PT framework, cannot exist independently of the scalar
potential $S$ from which they are derived.

Position C becomes thus, under PT, more than an interpretation:
it is the only framework in which the four natural conditions
(i)--(iv) of Proposition~\ref{prop:theta_p_unique_en} can be
satisfied simultaneously. The foundational debate is, in this
precise measure, adjudicated.

\section{CRT decoupling and modular non-locality}
\label{sec:crt_decoupling}

Sections~\ref{sec:field_connection} and~\ref{sec:holonomy_primitive}
have established the status of the holonomy angle $\theta_p$ as a
primitive invariant and that of the classical objects $A_\mu,
F_{\mu\nu}, g_{\mu\nu}$ as derivatives of the scalar potential
$S$. To connect this status to the experimental question of
non-locality, it remains to examine the correlations between
distinct modular channels --- the only ``inter-object''
correlations that PT allows to carry a non-trivial informational
significance. This section states, without detailed proof, the
two structural theorems proved in the technical companion
paper~\cite{PT_CRT_DECOUPLING}, together with their empirical
corollary.

\subsection{Theorem 4.1: Geometric Decoupling Theorem (empirical form)}
\label{ssec:thm_4_1_en}

\begin{theorem}[Geometric Decoupling Theorem, empirical form]
\label{thm:invariance_geom_en}
\footnote{Formally verified in Lean~4 + Mathlib as
\texttt{PT.CrtDecoupling.Empirical.empirical\_invariance}
(hypothesis form),
\texttt{SpectralReduction.empirical\_invariance\_closed}
(finite-state closed form),
\texttt{Phase3.empirical\_invariance\_infinite}
(infinite-state form via Birkhoff--Hopf), and
\texttt{Phase4.sieve\_empirical\_invariance}
(sieve-specific closure);
see \cite[\S~8.5]{PT_CRT_DECOUPLING}.}
Let $p, q \in \{3, 5, 7\}$ be distinct active primes and
$\Phi_p \in \mathcal{A}_p$, $\Phi_q \in \mathcal{A}_q$ bounded
observables on the corresponding modular $\sigma$-algebras. The
function
\begin{equation}
  \mathcal{I}_{p,q}(\mu)
  \;:=\;
  I_{\pi_m^{\mathrm{emp}}}\!\bigl(\Phi_p\,;\,\Phi_q\,\big|\,G(\mu)\bigr)
\end{equation}
is constant on $\mu \in [\mu_c, \infty)$.
\end{theorem}

The full proof is given as Theorem~6.1 of the technical
companion paper~\cite{PT_CRT_DECOUPLING}. It proceeds in four
steps: (i) the empirical measure $\pi_m^{\rm emp}$ of the residue
trajectory is a functional of the sole field $\{g_n\}$, hence
independent of $\mu$; (ii) the geometry $G(\mu)$ is measurable
with respect to the $\sigma$-algebra generated by the spectrum of
the transfer matrix $T_m$, which is itself contained in the
shift-invariant $\sigma$-algebra; (iii) the sieve dynamics is
ergodic by Birkhoff--Hopf applied to the spectral gap T4, hence
the invariant $\sigma$-algebra is trivial modulo
$\pi_m^{\rm emp}$, hence $G(\mu)$ is almost surely constant;
(iv) conditioning on an almost surely constant random variable
does not change mutual information. The conclusion follows.

The conceptual point is that the emergent metric $G(\mu)$ ---
including across the Hartle--Hawking transition at $\mu_c \approx
6.97$ and up to the canonical fixed point $\mustar = 15$ --- never
modifies the joint informational content of the active modular
channels. Inter-channel correlations are \emph{intrinsically}
pre-geometric: they exist before the construction of the metric
and are indifferent to it.

\subsection{Theorem 4.2: Geometric Decoupling Theorem (Ruelle form)}
\label{ssec:thm_4_2_en}

\begin{theorem}[Geometric Decoupling Theorem, Ruelle form]
\label{thm:ruelle_zero_en}
\footnote{Formally verified in Lean~4 + Mathlib as
\texttt{PT.CrtDecoupling.Main.geometric\_decoupling\_ruelle};
see \cite[\S~8.5]{PT_CRT_DECOUPLING}.}
Under the idealised measure $\pi_m^{\mathrm{Ruelle}}$, we have
$I(\Phi_p\,;\,\Phi_q) = 0$ exactly.
\end{theorem}

The proof is immediate from the CRT tensor factorisation of the
transfer $T_{pq} = T_p \otimes T_q$ (Theorem~3.1
of~\cite{PT_CRT_DECOUPLING}) and the corresponding factorisation
of the unique left Perron eigenvector
$\pi_{pq}^{\rm Ruelle} = \pi_p^{\rm Ruelle} \otimes \pi_q^{\rm Ruelle}$
(Lemma~3.2 of the same paper). This factorisation is an
independence condition, hence the Kullback--Leibler divergence
between the joint law and the product of marginals is zero; by
the data-processing inequality, the mutual information of any
pair of bounded observables inherits the vanishing.

\emph{Numerical validation.} The factorisation
$\pi_{15}^{\rm Ruelle} = \pi_3^{\rm Ruelle} \otimes \pi_5^{\rm Ruelle}$
is verified at the residual level
\begin{equation}
  \|\pi_{15}^{\rm direct} - \pi_3 \otimes \pi_5\|_2
  \;=\; 5.88 \times 10^{-53}
\end{equation}
at \texttt{mpmath} 50 digits of precision (test T2 PASS of
\texttt{scripts/test\_crt\_decoupling.py}, see
\cite{PT_CRT_DECOUPLING}).

\subsection{Corollary 4.3: Decoupling Closed Form}
\label{ssec:cor_4_3_en}

\begin{corollary}[Decoupling Closed Form]
\label{cor:closed_form_en}
\footnote{Not formalisable as a Lean theorem in the present form
because the depth parameter $k_{p,q}^{\rm eff}$ is fitted
empirically; only the prediction at a fixed $k$ admits a
Lean-verifiable computational statement;
see \cite[\S~8.5]{PT_CRT_DECOUPLING}.}
The constant $\mathcal{I}_{p,q}$ admits a closed-form expression
derived from the holonomy angles:
\begin{equation}
  \mathcal{I}_{p,q}
  \;=\;
  \frac{1}{2 \ln 2}\Bigl[
    \cos^{2 k_{p,q}}(\theta_3, \mu^\ast)
    + \sin^{4 k_{p,q}}(\theta_2, p_1)
  \Bigr].
  \label{eq:closed_form_en}
\end{equation}
Numerical validations (cf.\ \cite{PT_CRT_DECOUPLING}, Table~B):
$\mathcal{I}_{3,5} = 0.132$ bits (target $0.138$, error $4.5\%$);
$\mathcal{I}_{3,7} = 0.306$ bits (target $0.283$, error $8.1\%$).
\end{corollary}

The depth parameter $k_{p,q}$ is fitted empirically at
$k_{3,5}^{\rm eff} = 7$ and $k_{3,7}^{\rm eff} = 4$; a
combinatorial derivation from the antidiagonal structure of $T_3$
remains open (cf.\ Remark~7.1 of \cite{PT_CRT_DECOUPLING}). This
uncertainty affects only the \emph{quantitative value} of
$\mathcal{I}_{p,q}$; the structural result
(Theorem~\ref{thm:invariance_geom_en}) is independent of the
choice of $k_{p,q}$.

\subsection{Spatial non-locality dissolved into modular non-locality}
\label{ssec:dissolution_en}

The three preceding statements allow one to reformulate, in PT
terms, the status of the apparent ``non-locality'' of quantum
correlations. One can summarise the situation in three propositions
strictly parallel to the three positions of the foundational
debate (\S~\ref{sec:debate}), but with a different status:

\paragraph{(i) Inter-channel correlations exist.}
$\mathcal{I}_{3,5}$ and $\mathcal{I}_{3,7}$ are non-zero and of the
magnitude predicted by Corollary~\ref{cor:closed_form_en}. PT does
not deny the existence of EPR/Bell correlations; on the contrary,
it provides a quantitative prediction of them.

\paragraph{(ii) These correlations are of arithmetic origin.}
The closed form of Corollary~\ref{cor:closed_form_en} involves
only the holonomy angles $\theta_3, \theta_2$ evaluated at the
canonical scales $\mustar = 15$ and $p_1 = 3$. No spatial
integral, no metric, no distance between points of space enters
the expression. The observed non-locality is not a non-locality
in space: it is an intra-CRT correlation between different
modular channels.

\paragraph{(iii) These correlations are invariant under $G(\mu)$.}
Theorem~\ref{thm:invariance_geom_en} forbids the emergent
geometry from \emph{coupling} the modular channels. Whatever the
position in the geometric family $\mu \in [\mu_c, \infty)$ ---
including at any distance between detectors in a Bell
configuration --- the mutual information between distinct
channels remains the same.

These three propositions together dissolve spatial non-locality
in the strong sense of position B: quantum correlations do not
cross space because they are not in space. They are in the CRT
structure, which manifests in the emergent geometry $G$ by
projection without crossing it or using it as a causal carrier.
The apparent paradox of information ``faster than $c$'' in Bell
tests is an artefact of the spatial reading of a correlation that,
in the PT substrate, is not spatial. We return to the experimental
consequences of this dissolution in
\S~\ref{ssec:bell_consequence_en}.

\section{Falsifiable prediction QH1a}
\label{sec:qh1a_prediction}

Sections~\ref{sec:field_connection}--\ref{sec:crt_decoupling}
establish PT's structural position with respect to the
foundational debate: the holonomy angle is primitive, classical
objects are derivative, correlations between modular channels are
arithmetic and invariant under the emergent geometry. This
position is not merely interpretive: it entails a precise
experimental prediction, falsifiable on several existing
platforms, and structurally distinguishable from any prediction
of standard physics on a curved manifold.

\subsection{Statement of the prediction}
\label{ssec:qh1a_statement_en}

\begin{prediction}[QH1a: geometric invariance of MI]
\label{pred:qh1a_en}
In a multi-electronic atomic system probed by CRT spectroscopy
(channels $p = 3, 5, 7$), the mutual information
$I(n^{(p)}\,;\,n^{(q)})$ between modular occupation numbers:
\begin{enumerate}[label=(\roman*),nosep]
  \item is non-zero and given by the closed-form expression of
  Corollary~\ref{cor:closed_form_en} (arithmetic test);
  \item is invariant under any geometric transformation of the
  system (nuclear displacement, lattice deformation, detector
  separation) (geometric test).
\end{enumerate}
\end{prediction}

Condition (i) is the quantitative content of
Corollary~\ref{cor:closed_form_en}: the MI takes the values
$\mathcal{I}_{3,5} = 0.138$ bits and
$\mathcal{I}_{3,7} = 0.283$ bits (corpus values), predictable to
$\pm 20\%$ by the closed form with a single integer parameter
$k_{p,q}$ per pair. Any measurement failing to reproduce these
values (within the statistical precision of the experiment)
invalidates the prediction.

Condition (ii) is the structural content of
Theorem~\ref{thm:invariance_geom_en}: the MI must remain constant
under any reparametrisation of the spatial geometry. Any
measurement detecting a dependence of the MI on detector
separation, internuclear distance, or any other geometric
parameter, invalidates the prediction.

\subsection{Falsification criterion}
\label{ssec:qh1a_falsification_en}

Let $G_0, G_1$ be two distinct geometric configurations of the
same quantum system, and let
$\mathcal{I}_{p,q}(G_0), \mathcal{I}_{p,q}(G_1)$ be the measured
mutual information under each. Define
\begin{equation}
  \Delta\mathcal{I}_{p,q}
  \;:=\;
  \mathcal{I}_{p,q}(G_1) - \mathcal{I}_{p,q}(G_0)
  \label{eq:Delta_I_def_en}
\end{equation}
and let $\epsilon_{\rm stat}$ be the statistical uncertainty of
$\Delta\mathcal{I}_{p,q}$ as evaluated by the measurement
procedure. PT predicts
\begin{equation}
  \bigl|\Delta\mathcal{I}_{p,q}\bigr| \;<\; \epsilon_{\rm stat}
  \qquad \text{for every pair }(G_0, G_1).
  \label{eq:QH1a_predict_en}
\end{equation}
\emph{Falsification criterion.} If an experiment measures
$|\Delta\mathcal{I}_{p,q}| > 5 \epsilon_{\rm stat}$ for a pair of
distinct geometric configurations, prediction QH1a is falsified,
and with it the invariance
theorem~\ref{thm:invariance_geom_en}. Such a falsification would
also be a direct refutation of PT, since $\theta_p$ and
$\mathcal{A}_p$ are structural objects of the theory and cannot
be adjusted to absorb a geometric signal.

\subsection{Candidate experimental platforms}
\label{ssec:qh1a_platforms_en}

Three contemporary platforms offer access to the modular
occupation numbers $n^{(p)}$ with independent control of the
spatial geometry.

\paragraph{(a) Cold atoms in optical lattices.}
An optical lattice with variable spacing allows one to scan the
distance between adjacent potential wells while measuring the
modular occupation distribution of internal levels by fluorescence
spectroscopy. Geometry is controlled by the wavelength of the
trapping laser; the occupation numbers $n^{(p)}$ are accessible
via resonant measurement at the modular residue level. Expected
precision: $\epsilon_{\rm stat} \sim 10^{-3}$ bits on
$\mathcal{I}_{3,5}$ for $N \sim 10^4$ trapped atoms, ample
sensitivity to falsify QH1a at $5\sigma$ in the presence of
geometric coupling.

\paragraph{(b) Trapped ions at controlled separations.}
A chain of Yb$^+$ or Be$^+$ ions allows one to vary the inter-ion
separation with micrometre precision while measuring inter-channel
correlations by state tomography. The hyperfine structure of
internal levels provides the modular occupation numbers for
$p = 3$ (effective SU(3) colour channel) and $p = 5$ (flavour
channel in configurations with several valence electrons).
Expected precision: $\epsilon_{\rm stat} \sim 10^{-4}$ bits, the
tightest constraint of the three.

\paragraph{(c) Diatomic molecules at adjustable internuclear distance.}
A diatomic molecule probed by laser spectroscopy at sub-MHz
resolution allows one to vary the internuclear distance $R$ via
an external electric field (Stark effect) and to measure
simultaneously the modular occupation numbers of the electronic
levels. The signature sought is the absence of dependence of
$\mathcal{I}_{3,5}$ on $R$ over a range of several angstroms.
Expected precision: $\epsilon_{\rm stat} \sim 5 \times 10^{-3}$
bits, more modest but structurally complementary (the geometric
coupling expected of a standard theory is here dominated by the
vibrational-potential variation).

\subsection{Contrast with standard physics}
\label{ssec:qh1a_contrast_en}

A standard quantum theory on curved manifold predicts generically
$|d\mathcal{I}_{p,q}/dG| \neq 0$. The structural reason is that
matter fields are, in the standard formalism, sections of bundles
above the curved base, and their correlations inherit the metric
structure of the base via minimal gravitational couplings. The
effect is typically weak (suppressed by $1/M_{\rm Pl}^2$ or by
the local gravitational coupling), but it is \emph{non-zero} and
computable.

PT, by contrast, predicts $d\mathcal{I}_{p,q}/dG = 0$ identically.
The reason is that the modular channels $\mathcal{A}_p$ are
defined before the metric $G$ and are measurably anterior to it.
No function of the metric appears in the closed
form~\eqref{eq:closed_form_en}; the derivative with respect to $G$
is exactly zero.

This contrast makes QH1a a strong discriminating test between the
two frameworks. A positive measurement ($|d\mathcal{I}/dG| >
5\epsilon$) falsifies PT and confirms the standard framework; a
null measurement ($|d\mathcal{I}/dG| < \epsilon$) confirms PT at
the structural level and excludes any minimal geometric coupling
above the achieved precision. None of the current platforms has
yet performed the complete measurement --- in particular, none
has scanned the geometric dependence of the MI over several
orders of magnitude of separation. This is the principal
experimental gap to be closed in order to subject PT to a
genuine falsificationist test.

\section{Consequences for Aharonov--Bohm, Bell, EPR}
\label{sec:consequences_bell_ab_epr_en}

The three preceding structural results --- the holonomy angle as
primitive (Proposition~\ref{prop:theta_p_unique_en}), geometric
decoupling (Theorem~\ref{thm:invariance_geom_en}), and the
arithmetic closed form of the correlations
(Corollary~\ref{cor:closed_form_en}) --- modify the reading of
three classical experiments of foundational physics: the
Aharonov--Bohm configuration, the Bell tests, and the original
EPR configuration. We examine each briefly.

\subsection{Aharonov--Bohm: the phase is $\theta_3$, not $A_\mu$}
\label{ssec:ab_consequence_en}

In the PT reading, the interferometric phase shift
$\Delta\varphi$ observed in an Aharonov--Bohm
configuration~\cite{AharonovBohm1959} is not
$(e/\hbar) \oint A \cdot d\ell$ but $\theta_3$, the holonomy
angle of the modular channel $p = 3$ associated with the flux
enclosed by the particle trajectory. The identification between
the two formulas is tautological in the presence of the DIC
dictionary that translates PT units into SI units; what
distinguishes the PT reading is not the numerical value of the
phase shift but the identification of the \emph{physical support}
of the phase shift.

This identification has two consequences. First, the AB phase
is a strictly topological invariant of the path $\gamma$: it
depends only on the homotopy class of $\gamma$ in the complement
of the flux region, and is insensitive to any continuous
deformation of $\gamma$ within that class. This topology is not
a property of the metric $G(\mu)$; it is anterior to the metric,
in the sense of Theorem~\ref{thm:invariance_geom_en}. Second,
the absence of field $F = 0$ in the interior of an ideal solenoid
is no longer a paradox: $F_{\mu\nu}$ is not the object that
carries the phase information. It is $\theta_3$ that carries it,
and $\theta_3$ is non-zero even where $F$ vanishes, because
$\theta_3$ is defined by the modular connection of PT, which is
not itself a local function of the field.

\subsection{EPR/Bell: intra-channel correlations, not inter-channel}
\label{ssec:bell_consequence_en}

The Bell tests~\cite{Bell1964,Aspect1982,Hensen2015} establish
that the correlations observed between spatially separated
detectors exceed the CHSH bound~\cite{CHSH1969} expected of a
local hidden-variable theory. In the standard reading these
correlations are interpreted as the signature of a fundamental
non-locality of the quantum field: the collapse of an EPR state
upon left-side measurement instantaneously determines the
correlated right-side state, at any distance, and apparently
faster than light (but without allowing superluminal
communication, in accordance with no-signalling).

In the PT reading, this interpretation reverses. The two
measurements performed in a Bell test do \emph{not} act on two
distinct modular channels (in which case, by
Theorem~\ref{thm:invariance_geom_en}, their MI would be
geometry-invariant and of magnitude predicted by the closed
form~\eqref{cor:closed_form_en}). They act on the \emph{same}
modular channel --- typically $p = 2$ for polarisation/spin
experiments, or $p = 3$ for quark-colour experiments
(analogues). An intra-channel correlation is not the object of
Theorem~\ref{thm:invariance_geom_en}: it is a higher-order
correlation in the CRT structure, independent of any notion of
spatial separation, because the modular channel exists before
the construction of space by projection.

This reading is consistent with the empirical signature of the
most precise tests: the correlation observed between two
detectors separated by several kilometres~\cite{Hensen2015} is
strictly identical to that observed at a few metres
\cite{Aspect1982}, to within measurement precision. No significant
dependence on distance has ever been detected. This absence of
dependence is, in the PT reading, the structural expectation:
what is measured is not in space, so spatial separation cannot
modify the result.

\subsection{Dissolution of spatial non-locality}
\label{ssec:dissolution_consequence_en}

The ontological consequence of the two preceding subsections is
the \emph{dissolution} ([DISS] status) of strong-sense spatial
non-locality. The ``paradox'' of information apparently
transmitted faster than $c$ disappears because the correlations
do not propagate in space: they exist in the CRT structure, and
space is merely an emergent projection of that structure via the
metric $G(\mu)$.

The [DISS] status is, in the PT nomenclature, weaker than [THM]
but stronger than mere reinterpretation. A dissolution asserts
that the problem, as posed in the classical framework, is
ill-posed: its implicit premisses --- the existence of a
classical background space and the requirement that every
observed correlation propagate within it --- are not valid in
PT. The problem does not resolve in the sense that one would
have found the ``true'' classical explanation of the EPR
correlations; it dissolves in the sense that the correlations
are seen as objects that never belonged to the spatial domain.

The experimental horizon is then clear: PT predicts no signal
that could separate two spatially separated detectors beyond the
arithmetic CRT correlation, and categorically excludes any
signal that would depend on the metric $G(\mu)$ between the two
detectors. Any positive observation of such a dependence
falsifies PT and the reading presented here simultaneously.
Conversely, experimental confirmation of the geometric invariance
of EPR correlations --- i.e.\ prediction QH1a applied to the
relativistic regime --- would be a direct validation of the PT
reading of the foundational debate.

\section{Three ontological readings (R1, R2, R3)}
\label{sec:ontological_readings_en}

PT, by its nature as a purely informational theory derived from
a single arithmetic parameter $s = 1/2$, admits three ontological
readings compatible with its formal structure. The trichotomy is
laid out in detail in chapters~24 and~26 of the
monograph~\cite{PT_Monograph}; we recall it here briefly and
apply it specifically to the Aharonov--Bohm / Bell debate
developed in this article. The crucial point to emphasise --- and
the substance of~\S~\ref{ssec:reading_independence_en} --- is
that the proofs of
Sections~\ref{sec:field_connection}--\ref{sec:qh1a_prediction}
are logically independent of the choice of reading. The
trichotomy bears on the \emph{interpretation} of what the
theorems say, not on their validity.

\subsection{R1: PT as mathematical coincidence}
\label{ssec:R1_en}

Reading R1 holds that PT reproduces the correct empirical content
--- the Aharonov--Bohm phase, the EPR correlations, the absence
of geometric dependence of the MI --- by a structural coincidence
between the combinatorics of the Eratosthenes sieve and the
phenomenological structure of quantum electrodynamics. No
ontological commitment is made on the physical nature of the
angles $\theta_p$: they are mathematical quantities useful for
prediction, and nothing more. The theorems of the preceding
sections retain their validity --- they are mathematical
assertions about the structure of the sieve --- but their
correspondence with observed phenomena is treated as an
unexplained empirical fact.

This reading is the epistemologically most cautious: it commits
only the minimum required to account for the observations. Its
weakness is that it offers no \emph{explanation} of the
coincidence: why does the combinatorics of the sieve yield
exactly the right phenomenological structure? R1 does not answer
this question.

\subsection{R2: structuralism}
\label{ssec:R2_en}

Reading R2 holds that the Eratosthenes sieve and quantum
electrodynamics share an abstract holonomic \emph{pattern}, and
that neither is ontologically prior. The $\theta_p$ are then
neither mathematical objects (in the sense of R1) nor physical
objects (in the sense of R3): they are components of an abstract
structure of which the two theories are incarnations.

This reading, classical in philosophy of science under the name
\emph{structuralism}~\cite{Belot1998,Maudlin2011}, finds in PT a
particularly favourable ground: the correspondence between the
sieve and the phenomenology is no longer a coincidence but the
expression of a common underlying structure, and the primacy of
holonomy becomes a structural property rather than a matter of
material ontology. The theorems of the preceding sections retain
their validity and acquire explanatory status in the sense that
they describe the common structure.

\subsection{R3: informational realism}
\label{ssec:R3_en}

Reading R3 holds that the $\theta_p$ \emph{are} the fundamental
physical reality, and that the classical objects of
electrodynamics --- $A_\mu$, $F_{\mu\nu}$, $g_{\mu\nu}$ --- are
phenomenological projections onto the emergent spacetime. Under
this reading, the ontological substance of the theory is
information ($\DKL$ in bits); the notions of particle, field,
space, and time are emergent constructions above this
informational substance.

R3 is the most committed reading: it asserts that the world is
made of arithmetic information, that matter and space are mere
projections of it, and that the correlations observed in Bell
tests are direct reflections of the fundamental CRT structure.
Under R3, the dissolution of spatial non-locality
(\S~\ref{ssec:dissolution_consequence_en}) takes its strongest
sense: the EPR correlations are literally non-spatial, because
what is measured has no spatial support.

R3 inherits the classical philosophical difficulties of
informational realism~\cite{Healey2007}: the nature of the
``execution'' of the sieve remains open (who or what
``computes'' the sequence $\{g_n\}$?), and the apparent
correspondence between the mathematical structure and the
observed phenomenology demands, under R3, a metaphysical
explanation that exceeds the scope of the theory itself.

\subsection{Theorem independence from the reading}
\label{ssec:reading_independence_en}

The central logical point of this section is that
Theorems~\ref{thm:invariance_geom_en}
and~\ref{thm:ruelle_zero_en},
Corollary~\ref{cor:closed_form_en}, and
Prediction~\ref{pred:qh1a_en} are \emph{all} valid under R1, R2,
R3. No formal result of the preceding sections depends on the
choice of ontological interpretation. The proofs are mathematical;
the predictions are functions of holonomy angles and the
emergent metric; and the numerical values (e.g.\
$\mathcal{I}_{3,5} = 0.138$ bits) are computable independently of
the chosen reading.

What changes with the reading is solely the \emph{interpretation}
of what the experiment measures. Under R1, an experimental
confirmation of QH1a is an unexplained empirical fact of
coincidence with the sieve structure. Under R2, it is the
empirical expression of an abstract common structure between
sieve and electrodynamics. Under R3, it is a direct measurement
of the ontological informational reality. The three readings are
compatible with the same observations; the choice among them
depends on extra-empirical criteria (parsimony, explanatory power,
metaphysical commitment).

This independence has an important practical consequence for the
experimental evaluation of PT: a test of QH1a can be conducted,
and its results interpreted, without prior commitment on the
ontological status of the $\theta_p$. The theory is testable
independently of the reading. To our knowledge, this is the
first theoretical framework that proposes a structural
dissolution of the Aharonov--Bohm / Bell debate while remaining
empirically falsifiable, and that can be evaluated without first
having to choose between the classical ontological positions.


\bibliographystyle{plain}
\bibliography{references}

\end{document}
